
\documentclass{article}
\usepackage{colortbl}
\usepackage{makecell}
\usepackage{multirow}
\usepackage{supertabular}

\begin{document}

\newcounter{utterance}

\centering \large Interaction Transcript for game `hot\_air\_balloon', experiment `air\_balloon\_survival\_en\_complexity\_easy', episode 0 with Qwen3{-}32B{-}t0.0.
\vspace{24pt}

{ \footnotesize  \setcounter{utterance}{1}
\setlength{\tabcolsep}{0pt}
\begin{supertabular}{c@{$\;$}|p{.15\linewidth}@{}p{.15\linewidth}p{.15\linewidth}p{.15\linewidth}p{.15\linewidth}p{.15\linewidth}}
    \# & $\;$A & \multicolumn{4}{c}{Game Master} & $\;\:$B\\
    \hline

    \theutterance \stepcounter{utterance}  
    & & \multicolumn{4}{p{0.6\linewidth}}{
        \cellcolor[rgb]{0.9,0.9,0.9}{
            \makecell[{{p{\linewidth}}}]{
                \texttt{\tiny{[P1$\langle$GM]}}
                \texttt{You are participating in a collaborative negotiation game.} \\
\\ 
\texttt{Together with another participant, you must agree on a single set of items that will be kept. Each of you has your own view of how much each item matters to you (importance). You do not know how the other participant values the items. Additionally, you are given the effort each item demands.} \\
\texttt{You may only agree on a set if the total effort of the selected items does not exceed a shared limit:} \\
\\ 
\texttt{LIMIT = 7605} \\
\\ 
\texttt{Here are the individual item effort values:} \\
\\ 
\texttt{Item effort = \{"B21": 369, "B00": 951, "A08": 156, "A27": 266, "B94": 581, "B53": 287, "A23": 866, "C65": 180, "A89": 800, "B67": 739, "B45": 637, "C34": 86, "C12": 749, "A05": 370, "B60": 345, "C92": 951, "B69": 147, "A84": 265, "C06": 262, "B12": 259, "A16": 358, "B85": 394, "C35": 286, "B10": 580, "C15": 479, "B09": 14, "B16": 153, "C81": 977, "A58": 134, "C24": 982, "A30": 259, "A66": 232, "B75": 202, "B46": 73, "C70": 821\}} \\
\\ 
\texttt{Here is your personal view on the importance of each item:} \\
\\ 
\texttt{Item importance values = \{"B21": 138, "B00": 583, "A08": 868, "A27": 822, "B94": 783, "B53": 65, "A23": 262, "C65": 121, "A89": 508, "B67": 780, "B45": 461, "C34": 484, "C12": 668, "A05": 389, "B60": 808, "C92": 215, "B69": 97, "A84": 500, "C06": 30, "B12": 915, "A16": 856, "B85": 400, "C35": 444, "B10": 623, "C15": 781, "B09": 786, "B16": 3, "C81": 713, "A58": 457, "C24": 273, "A30": 739, "A66": 822, "B75": 235, "B46": 606, "C70": 968\}} \\
\\ 
\texttt{Goal:} \\
\\ 
\texttt{Your goal is to negotiate a shared set of items that benefits you as much as possible (i.e., maximizes total importance to YOU), while staying within the LIMIT. You are not required to make a PROPOSAL in every message {-} you can simply negotiate as well. All tactics are allowed!} \\
\\ 
\texttt{Interaction Protocol:} \\
\\ 
\texttt{You may only use the following structured formats in a message:} \\
\\ 
\texttt{PROPOSAL: \{'A', 'B', 'C', …\}} \\
\texttt{Propose keeping exactly those items.} \\
\texttt{REFUSE: \{'A', 'B', 'C', …\}} \\
\texttt{Explicitly reject opponent's proposal.} \\
\texttt{ARGUMENT: \{'...'\}} \\
\texttt{Defend your last proposal or argue against the player's proposal.} \\
\texttt{AGREE: \{'A', 'B', 'C', …\}} \\
\texttt{Accept the opponent's proposal which ends the game.} \\
\texttt{STRATEGIC REASONING: \{'...'\}} \\
\texttt{	Describe your strategic reasoning or anticipation explaining your choice of action. This is a hidden message which will not be shared with the other participant.} \\
\\ 
\texttt{Rules:} \\
\\ 
\texttt{You may only AGREE on a proposal the other party has logged via PROPOSAL.} \\
\texttt{You may only REFUSE a proposal the other party has logged via PROPOSAL.} \\
\texttt{Total effort of any PROPOSAL or AGREE set must be ≤ LIMIT.} \\
\texttt{Do NOT reveal your hidden importance scores.} \\
\texttt{A tag in a structured format must be followed by colon and whitespace. The argument must be a python set containing 0 or more strings.} \\
\texttt{So, it must be of the form TAG: \{...\}} \\
\texttt{Strictly follow the interaction protocol and DO NOT write anything beyond the given structure.} \\
\texttt{The game ends when one side gives an AGREE to a PROPOSAL made by the other player.} \\
\texttt{The content in your response which can be handed to the other player has to be non{-}empty.} \\
\texttt{Only proposals which have been logged via the PROPOSAL format structure and which haven't been refused via REFUSE are active.} \\
\texttt{You must include the ARGUMENT format at least once somewhere in all of your messages.} \\
\texttt{You must include the STRATEGIC REASONING format only once at the very beginning of every one of your messages and not more often. The contents will not be given to the other player so they can include anything you like including your own importance values. Here you should reason step by step to come up with you next move.} \\
            }
        }
    }
    & \\ \\

    \theutterance \stepcounter{utterance}  
    & \multicolumn{4}{p{0.6\linewidth}}{
        \cellcolor[rgb]{0.9,0.9,0.9}{
            \makecell[{{p{\linewidth}}}]{
                \texttt{\tiny{[P1$\rangle$GM]}}
                \texttt{<think>} \\
\texttt{Okay, let's see. I need to figure out the best set of items to propose that maximizes my total importance without exceeding the effort limit of 7605. First, I should calculate the total effort of all items to see if the sum is over the limit. Let me add them up.} \\
\\ 
\texttt{Looking at the item efforts: B21 is 369, B00 is 951, A08 is 156, A27 is 266, B94 is 581, B53 is 287, A23 is 866, C65 is 180, A89 is 800, B67 is 739, B45 is 637, C34 is 86, C12 is 749, A05 is 370, B60 is 345, C92 is 951, B69 is 147, A84 is 265, C06 is 262, B12 is 259, A16 is 358, B85 is 394, C35 is 286, B10 is 580, C15 is 479, B09 is 14, B16 is 153, C81 is 977, A58 is 134, C24 is 982, A30 is 259, A66 is 232, B75 is 202, B46 is 73, C70 is 821.} \\
\\ 
\texttt{Adding all these up... Let me do this step by step. Maybe group them to make it easier. Let's see:} \\
\\ 
\texttt{Start with B21 (369) + B00 (951) = 1320. Then A08 (156) + A27 (266) = 422. Total so far: 1742. B94 (581) + B53 (287) = 868. Total: 2610. A} \\
            }
        }
    }
    & & \\ \\

    \theutterance \stepcounter{utterance}  
    & & \multicolumn{4}{p{0.6\linewidth}}{
        \cellcolor[rgb]{0.9,0.9,0.9}{
            \makecell[{{p{\linewidth}}}]{
                \texttt{\tiny{[P1$\langle$GM]}}
                \texttt{Your response did not start with the proper strategic reasoning tag at the very beginning of your response. The very first structured format must be of the form STRATEGIC REASONING: \{...\}. Try again.} \\
            }
        }
    }
    & \\ \\

    \theutterance \stepcounter{utterance}  
    & \multicolumn{4}{p{0.6\linewidth}}{
        \cellcolor[rgb]{0.9,0.9,0.9}{
            \makecell[{{p{\linewidth}}}]{
                \texttt{\tiny{[P1$\rangle$GM]}}
                \texttt{<think>} \\
\texttt{Okay, let me start over. The user wants me to make sure that my response begins with the STRATEGIC REASONING tag. I need to outline my thought process first, then proceed with the action.} \\
\\ 
\texttt{First, I need to calculate the total effort of all items to see if the sum is over the limit. If it is, I can't take all. Let me add up all the effort values. } \\
\\ 
\texttt{Looking at the item efforts: B21 is 369, B00 is 951, A08 is 156, A27 is 266, B94 is 581, B53 is 287, A23 is 866, C65 is 180, A89 is 800, B67 is 739, B45 is 637, C34 is 86, C12 is 749, A05 is 370, B60 is 345, C92 is 951, B69 is 147, A84 is 265, C06 is 262, B12 is 259, A16 is 358, B85 is 394, C35 is 286, B10 is 580, C15 is 479, B09 is 14, B16 is 153, C81 is 977, A58 is 134, C24 is 982, A30 is 259, A66 is 232, B75 is 202, B46 is 73, C70 is 821.} \\
\\ 
\texttt{Adding them up step by step. Let's group them to make it easier. } \\
\\ 
\texttt{Start with B21 (369) + B00 (951) = 1320. Then A08 (156) + A27 (266) = 422. Total so far: 1742. B94 (581) + B53 (287) =} \\
            }
        }
    }
    & & \\ \\

    \theutterance \stepcounter{utterance}  
    & & \multicolumn{4}{p{0.6\linewidth}}{
        \cellcolor[rgb]{0.9,0.9,0.9}{
            \makecell[{{p{\linewidth}}}]{
                \texttt{\tiny{[P1$\langle$GM]}}
                \texttt{Your response did not start with the proper strategic reasoning tag at the very beginning of your response. The very first structured format must be of the form STRATEGIC REASONING: \{...\}. Try again.} \\
            }
        }
    }
    & \\ \\

    \theutterance \stepcounter{utterance}  
    & \multicolumn{4}{p{0.6\linewidth}}{
        \cellcolor[rgb]{0.9,0.9,0.9}{
            \makecell[{{p{\linewidth}}}]{
                \texttt{\tiny{[P1$\rangle$GM]}}
                \texttt{<think>} \\
\texttt{Okay, I need to make sure I start with the STRATEGIC REASONING tag. Let me outline my plan step by step.} \\
\\ 
\texttt{First, I'll calculate the total effort of all items to see if the sum exceeds the limit. If it does, I'll need to find a subset that stays under 7605. Then, I'll prioritize items with high importance to me while keeping effort in check. Maybe start by selecting the items with the highest importance{-}to{-}effort ratio. Let me check the importance values and see which ones are the most valuable. For example, A08 has a high importance of 868 with low effort 156. Similarly, A27 has 822 and 266. I should include those. Then look for other high importance items that don't add too much effort. Need to avoid items with low importance even if they have low effort. Also, check for any items with very high effort that might be worth excluding. Once I have a tentative list, calculate the total effort and adjust as needed. Maybe start with the top 10 items by importance and see if the effort is under. If not, swap out the least important ones. Let's try that.} \\
\texttt{</think>} \\
\\ 
\texttt{STRATEGIC REASONING: \{  } \\
\texttt{  "I will calculate the total effort of all items to determine if the full set is feasible. If not, I'll prioritize items with the highest importance{-}to{-}effort ratio. My top priorities are A08 (importance 868, effort 156), A27 (822, 266), and A16 (856, 358). I'll also consider high{-}importance items like B09 (786, 14) and C70 (968, 821). I'll exclude low{-}importance items like B16 (3, 153) and C06 (30, 262). I'll test combinations to stay under 7605 while maximizing my total importance."  } \\
\texttt{\}  } \\
\texttt{PROPOSAL: \{"A08", "A27", "A16", "B09", "C70", "A05", "B60", "C15", "B67", "B94"\}  } \\
\texttt{ARGUMENT: \{"This proposal} \\
            }
        }
    }
    & & \\ \\

    \theutterance \stepcounter{utterance}  
    & & & \multicolumn{2}{p{0.3\linewidth}}{
        \cellcolor[rgb]{0.9,0.9,0.9}{
            \makecell[{{p{\linewidth}}}]{
                \texttt{\tiny{[GM$|$GM]}}
                \texttt{abort} \\
            }
        }
    }
    & & \\ \\

    \theutterance \stepcounter{utterance}  
    & & & \multicolumn{2}{p{0.3\linewidth}}{
        \cellcolor[rgb]{0.9,0.9,0.9}{
            \makecell[{{p{\linewidth}}}]{
                \texttt{\tiny{[GM$|$GM]}}
                \texttt{end game} \\
            }
        }
    }
    & & \\ \\

\end{supertabular}
}

\end{document}
